\scanpage{270}
\mspage{2}

\uline{Loc} 3 (\uline{localisation}, forme faible) : Soit $f:X\longrightarrow S$ une
flèche \unsure{dans|de} $\mathrm{Cat}$ $W$-asphérique, i.e. telle que
$\forall\,s\in\mathrm{Ob}\,S$, \inserted{la cat. localisée $X_{/s}$ de $X$ \unsure{sur|en} $s$} $X_{/s}=X\times_S S_{/s}$ soit $W$-asphérique. Alors $f\in W$.

% slanted
\margin{Dans une version finale, il faudrait intervertir les notations pour $\mathrm{Loc}(3)$ et pour $\mathrm{Loc}(3\,\mathrm{bis})$. Car ce dernier m'apparaît comme le plus important des deux -- il est \unsure{obligé|satisfait}, pour \unsure{obtenir|tous} les localiseurs \deleted{\unsure{qui}} \unsure{ayant|peuvent} \unsure{été|être} utilisés. \unsure{Ceci|Ce} au moins (\unsure{sauf erreur}), est \deleted{\unsure{va}}valable pour tt localiseur de la forme $W_{\unsure{\mathbb D|D}}$, où $\mathbb D$ est un dérivateur (ou seulement un prédérivateur satisfaisant $\mathbb D\unsure{\mathrm{er}}\,(4,5)$).}

\uline{Rappels}. 1. $\mathrm{Loc}(1,2,3)$ impliquent que l'on a

\[
W=W^o\qquad\text{i.e}\quad f\in W\iff f^o\in W,
\]

\unsure{dénotant|dénoté} par $X\longmapsto X^o$, $f\longmapsto f^o$ le
foncteur de passage à la catégorie opposée, \inserted{\uline{foncteur}} qui

\[
X\longmapsto X^o:\quad \mathrm{Cat}\xrightarrow{\sim}\mathrm{Cat}
\]

qui est un \uline{automorphisme} de $\mathrm{Cat}$ (donc
commute à toutes $\underrightarrow{\lim}$ et $\underrightarrow{\lim}$). Il
s'ensuit que l'ensemble d'axiomes
$\mathrm{Loc}(1,2,3)$ équivaut à $\mathrm{Loc}(1,2_{\unsure{0|o}}^{\mathrm{opp}},3^{\mathrm{opp}})$, où
$\mathrm{Loc}(2^{\mathrm{opp}})$ est l'axiome ``dual'' de $\mathrm{Loc}(2)$,
savoir que \deleted{tt \unsure{catégorie}} objet de $\mathrm{Cat}$ ayant
un objet initial est $W$-asphérique, et
\unsure{de m\textsuperscript{e}|de =}
$\mathrm{Loc}(3^{\mathrm{opp}})$ l'axiome que \deleted{\unsure{tout}} $f:X\to S$ \uline{$W$-coasphérique}, i.e telle que les \inserted{catégories colocalisées}
$ {}_{s\backslash}X$ \inserted{de $X$ sous $s$} \deleted{\unsure{sont}} \deleted{\unsure{soient}} $W$-asphériques, est dans $W$.

\deleted{\unsure{Loc (3 bis)}} \deleted{\unsure{mutuel}}
\uline{Définition} 1.1. On appelle \uline{localiseur fondamental}
une partie $W$ de $\mathrm{Fl}(\mathrm{Cat})$ satisfaisant les
axiomes $\mathrm{Loc}(1,2,3)$ (i.e. saturation modérée,
axiome de l'objet final, axiome de localisation).
% archive sheet number: 269
